% Generated from locale/tr/content/computability/recursive-functions/notation-pr-functions.tex; do not hand-edit.


\olfileid[tr]{cmp}{rec}{not}
\olsection{İlkel Özyinelemeli Fonksiyonlar İçin Gösterimler}

İlkel özyinelemeli fonksiyonların kesin bir tümevarımsal tanımına sahip
olmanın yararlarından biri, onları sistematik biçimde betimleyebilmemizdir.
Örneğin her böyle fonksiyona şu şekilde bir ``gösterim'' atayabiliriz. Sıfır,
ardıl ve izdüşümler için $\Zero$, $\Succ$ ve $\Proj{n}{i}$ simgelerini
kullanalım. Şimdi $h$ fonksiyonunun $k$-yerli $f$ ile $n$-yerli
$g_0,\dots,g_{k-1}$ fonksiyonlarından bileşkeyle tanımlandığını ve bu son
fonksiyonlara sırasıyla $F,G_0,\dots,G_{k-1}$ gösterimlerini atadığımızı
varsayalım. Yeni bir $\fn{Comp}_{k,n}$ simgesi kullanarak $h$ fonksiyonunu
$\fn{Comp}_{k,n}[F,G_0,\dots,G_{k-1}]$ ile gösterebiliriz.

İlkel özyinelemeyle tanımlanan fonksiyonlar için benzer gösterimler
kullanabiliriz. $(k+1)$-yerli $h$ fonksiyonu, $k$-yerli $f$ ile
$(k+2)$-yerli $g$ fonksiyonlarından ilkel özyinelemeyle tanımlanmış ve
$f$ ile $g$'ye sırasıyla $F$ ile $G$ gösterimleri atanmış olsun. O zaman
$h$'ye atanan gösterim $\fn{Rec}_k[F,G]$ olur.

Toplama fonksiyonunun ilkel özyinelemeyle şöyle tanımlandığını anımsayalım:
\begin{align*}
  \Add(x_0,0) & = \Proj{1}{0}(x_0)=x_0 \\
  \Add(x_0,y+1) & = \Succ(\Proj{3}{2}(x_0,y,\Add(x_0,y)))
    =\Add(x_0,y)+1.
\end{align*}
Burada $f$'nin rolünü $\Proj{1}{0}$, $g$'nin rolünü ise
$\Succ(\Proj{3}{2}(x_0,y,z))$ oynar. Sonuncusu, bir yerli $\Succ$ ile üç
yerli $\Proj{3}{2}$ fonksiyonlarından bileşkeyle tanımlandığı için
$\fn{Comp}_{1,3}[\Succ,\Proj{3}{2}]$ gösterimini alır. Bu düzenle toplama
fonksiyonunu
\[
  \fn{Rec}_1[\Proj{1}{0},\fn{Comp}_{1,3}[\Succ,\Proj{3}{2}]]
\]
ile gösterebiliriz. Bu gösterimler, örneğin ilkel özyinelemeli fonksiyonları
numaralandırırken yararlı olur.

\begin{prob}
  $\Mult$ için eksiksiz ilkel özyinelemeli gösterimi verin.
\end{prob}

